\clearpage
\VerdierPhysicalPageAnchor{0117}%
\begingroup
\small
\setlength{\abovedisplayskip}{4pt plus 2pt minus 2pt}%
\setlength{\belowdisplayskip}{4pt plus 2pt minus 2pt}%

\VerdierSemanticAnchor{ch2-1-2-5-combined-diagram}
On en déduit que le diagramme ci-après est commutatif :
\begin{equation*}
  \resizebox{0.98\linewidth}{!}{%
    $\begin{tikzcd}[ampersand replacement=\&,column sep=2.1em,row sep=3.1em]
      X\oplus X' \arrow[r,"{u\oplus u'}"] \arrow[d,"{\operatorname{id}}"']
        \& Y\oplus Y' \arrow[r,"{v\oplus v'}"] \arrow[d,"{\operatorname{id}}"']
        \& Z\oplus Z' \arrow[r,"{w\oplus w'}"] \arrow[d,"{(f,f')}"']
        \& (X\oplus X')[1] \arrow[d,"{\operatorname{id}}"] \\
      X\oplus X' \arrow[r,"{u\oplus u'}"']
        \& Y\oplus Y' \arrow[r,"m"']
        \& L \arrow[r,"n"']
        \& (X\oplus X')[1]
    \end{tikzcd}$%
  }\,.
\end{equation*}

\VerdierSemanticAnchor{ch2-1-2-5-sum-cone-map}
Pour démontrer le corollaire, il suffit de montrer que le morphisme :
\[
  (f,f')\colon Z\oplus Z'\longrightarrow L
\]
est un isomorphisme. Il suffit donc de montrer que, pour tout objet $M$, le
morphisme :

\VerdierSemanticAnchor{ch2-1-2-5-hom-test}
\[
  \operatorname{Hom}_{\mathcal D}\bigl(M,(f,f')\bigr)
\]
est un isomorphisme de groupes abéliens ; ceci résulte immédiatement de
(1.2.1) et du lemme des cinq.

\VerdierSemanticAnchor{ch2-1-2-5-proof-complete}

\VerdierSemanticAnchor{ch2-1-2-6}
\medskip
\noindent\textbf{Corollaire 1.2.6.}\quad
Soit
$X\xrightarrow{u}Y\xrightarrow{v}Z\xrightarrow{w}X[1]$
un triangle distingué de $\mathcal D$ tel que $w=0$. Il existe alors une
section $s\colon Z\longrightarrow Y$ du morphisme $v$. Toute section
$s\colon Z\longrightarrow Y$ définit un isomorphisme :

\VerdierSemanticAnchor{ch2-1-2-6-split-isomorphism}
\[
  X\oplus Z\xrightarrow[\sim]{(u,s)}Y\,.
\]

\VerdierSemanticAnchor{ch2-1-2-6-converse}
Réciproquement, tout triangle du type :

\VerdierSemanticAnchor{ch2-1-2-6-1}
\begin{equation}
  X\xrightarrow{\binom{\operatorname{id}}{0}}X\oplus Z
  \xrightarrow{(0,\operatorname{id})}Z\xrightarrow{0}X[1]
  \tag{1.2.6.1}
\end{equation}
est distingué.

\VerdierSemanticAnchor{ch2-1-2-6-third-assertion}
La troisième assertion résulte immédiatement de TRI, de TRII et de (1.2.5).

\VerdierSemanticAnchor{ch2-1-2-6-first-assertion}
Démontrons la première assertion. Il résulte de (1.2.1) que, pour tout $W$,
la suite :

\VerdierSemanticAnchor{ch2-1-2-6-2}
\begin{equation}
  0\longrightarrow\operatorname{Hom}_{\mathcal D}(W,X)
  \longrightarrow\operatorname{Hom}_{\mathcal D}(W,Y)
  \longrightarrow\operatorname{Hom}_{\mathcal D}(W,Z)
  \longrightarrow0
  \tag{1.2.6.2}
\end{equation}
\endgroup
